% Figura: los campos de referencia con un volumen V (cubo) y un circuito C.
\newcommand{\cuboV}[2]{%  esquina inferior izquierda (#1,#2): cubo de lado 1
  \fill[grisclaro] (#1,#2) rectangle (#1+1,#2+1);
  \fill[grispizarra!40] (#1,#2+1) -- (#1+0.35,#2+1.3) -- (#1+1.35,#2+1.3) -- (#1+1,#2+1) -- cycle;
  \fill[grispizarra!55] (#1+1,#2) -- (#1+1.35,#2+0.3) -- (#1+1.35,#2+1.3) -- (#1+1,#2+1) -- cycle;
  \draw[rojoteo, line width=0.8pt] (#1,#2) rectangle (#1+1,#2+1);
  \draw[rojoteo, line width=0.8pt] (#1,#2+1) -- (#1+0.35,#2+1.3) -- (#1+1.35,#2+1.3) -- (#1+1.35,#2+0.3) -- (#1+1,#2);
  \draw[rojoteo, line width=0.8pt] (#1+1,#2+1) -- (#1+1.35,#2+1.3);
  \node[font=\footnotesize, text=rojoteo] at (#1+0.5,#2+0.5) {$V$};}
\newcommand{\circuitoC}[4]{%  rectángulo [#1,#2]x[#3,#4] recorrido en sentido horario
  \fill[white, opacity=0.55] (#1,#3) rectangle (#2,#4);\draw[rojoteo, line width=1pt] (#1,#3) rectangle (#2,#4);
  \draw[-{Stealth[length=6pt]}, rojoteo, line width=1pt] ({(#1+#2)/2-0.15},#4) -- ({(#1+#2)/2+0.05},#4);
  \draw[-{Stealth[length=6pt]}, rojoteo, line width=1pt] (#2,{(#3+#4)/2+0.12}) -- (#2,{(#3+#4)/2-0.05});
  \draw[-{Stealth[length=6pt]}, rojoteo, line width=1pt] ({(#1+#2)/2+0.15},#3) -- ({(#1+#2)/2-0.05},#3);
  \draw[-{Stealth[length=6pt]}, rojoteo, line width=1pt] (#1,{(#3+#4)/2-0.12}) -- (#1,{(#3+#4)/2+0.05});
  \node[font=\footnotesize, text=rojoteo, anchor=west] at (#2+0.05,{(#3+#4)/2}) {$C$};}
\begin{figure}[ht]
\centering
\begin{minipage}[t]{0.325\textwidth}\centering
\begin{tikzpicture}[x=0.68cm, y=0.68cm, baseline=(current bounding box.north)]
  \marcoD{-2.7}{2.7}{-2.7}{2.7}{-2,-1,0,1,2}{-2,-1,0,1,2}
  \ejesD{-2.7}{2.7}{-2.7}{2.7}
  \foreach \x in {-2.4,-1.8,...,2.4}{\foreach \y in {-2.4,-1.8,...,2.4}{%
    \draw[-{Stealth[length=4.5pt,width=4.5pt]}, line width=0.7pt, azulprof]
      (\x-0.17,\y-0.17) -- ++(0.34,0.34);}}
\end{tikzpicture}\\[3pt]{\small $\vc{F}_1=\ux+\uy$}
\end{minipage}\hfill
\begin{minipage}[t]{0.325\textwidth}\centering
\begin{tikzpicture}[x=1.12cm, y=1.12cm, baseline=(current bounding box.north)]
  \marcoD{0}{3.2}{0}{3.2}{0,1,2,3}{0,1,2,3}
  \ejesD{0}{3.2}{0}{3.2}
  \clipD{0}{3.2}{0}{3.2}{\flechas{\x}{0}{0.2,0.575,...,3.1}{0.2,0.575,...,3.1}{0.30}{3.1}}
  \cuboV{1}{1.5}
  \circuitoC{1}{2}{0.5}{1}
\end{tikzpicture}\\[3pt]{\small $\vc{F}_2=x\,\ux$}
\end{minipage}\hfill
\begin{minipage}[t]{0.325\textwidth}\centering
\begin{tikzpicture}[x=1.12cm, y=1.12cm, baseline=(current bounding box.north)]
  \marcoD{0}{3.2}{0}{3.2}{0,1,2,3}{0,1,2,3}
  \ejesD{0}{3.2}{0}{3.2}
  \clipD{0}{3.2}{0}{3.2}{\flechas{\y}{0}{0.2,0.575,...,3.1}{0.2,0.575,...,3.1}{0.30}{3.1}}
  \cuboV{1}{1.5}
  \circuitoC{1}{2}{0.5}{1}
\end{tikzpicture}\\[3pt]{\small $\vc{F}_3=y\,\ux$}
\end{minipage}
\caption{Campos de referencia con un volumen $V$ (cubo) y un circuito $C$
(rectángulo recorrido en sentido horario). Para $\vc{F}_2$ el flujo neto por la
superficie del cubo es no nulo ($=1$) y la circulación por $C$ es nula; para
$\vc{F}_3$ ocurre lo contrario: el flujo es nulo y la circulación vale $0{,}5$.}
\label{fig:campos_int}
\end{figure}
